% ===========================================================================
%  sviluppatore/05-estendere/07-metrica.tex
% ===========================================================================
\section{Una metrica: i dati del sistema e la pagina dello stream}\label{sez:nuova-metrica}

Una metrica è il più piccolo degli artefatti: un file \file{metrics/<nome>.py}
con un \cod{META} di soli letterali e, sotto, codice vettoriale sull'equity
della pagina. Produce un numero (\cod{result(...)}) o uno o più grafici
(\cod{plot}, \cod{bars}, \cod{heatmap}, \cod{chart}). La si scrive in Studio,
che il \cod{META} lo compone da sé con i campi \voce{Output} e \voce{Chart
type}; il contratto completo, con gli esempi, è in
\file{metrics/METRIC\_TEMPLATE.md}.

Gira in \file{dashboard/metric\_runtime.py}, in una sandbox: builtins ridotti,
niente \cod{import}, niente attributi \cod{\_\_dunder\_\_} né privati. Fino a
questa versione vedeva soltanto l'equity della pagina e, quando la pagina lo
passava, il benchmark. Ora legge anche i dati del sistema.

\subsection{L'oggetto \cod{data}}

Nel codice di una metrica c'è \cod{data}, istanza di \cod{Data}
(\file{dashboard/metric\_data.py}). Ogni cartella di \pkgDati{} è una
categoria e ogni categoria è una classe, tutte con gli stessi metodi:

\begin{center}
\begin{tabular}{@{}L{36mm}L{102mm}@{}}
\toprule
\textbf{Metodo} & \textbf{Che cosa restituisce} \\
\midrule
\rigaTenue
\cod{names()} & i file della categoria: simboli per le categorie 1:1, file
  condivisi per le 1:N (\cod{settore\_tech}, \cod{USD}) \\
\cod{timeframes(nome)} & le sottocartelle di timeframe presenti, dalla più
  fine \\
\rigaTenue
\cod{columns(nome)} & i campi veri di un file \\
\cod{get(nome)} & il file com'è sul disco (con \cod{start} ed \cod{end} se
  ne ritaglia un periodo) \\
\rigaTenue
\cod{series(nome, campo)} & un campo, allineato alle date della metrica
  (\cod{align=False} lo dà grezzo) \\
\bottomrule
\end{tabular}
\end{center}

\noindent
Tranne \cod{timeframes}, accettano tutti \cod{timeframe=}: senza, vale quello
predefinito (sotto).
Le categorie del sistema hanno una classe propria: \cod{Assets} (con
\cod{close}, \cod{returns}, \cod{ohlcv}), \cod{Earnings} (con \cod{events}),
\cod{Fundamentals}, \cod{Rates}, \cod{Macro} (anche i campi del driver
macro), \cod{News} (titoli al giorno), \cod{MMS} (file \cod{<SIMBOLO>\_<TF>}).
Quelle create dalla pagina Dati sono \cod{Category}.

\begin{codiceInterno}{python}{Una metrica che legge un asset}{metrics/beta\_spy.py}
META = {"label": "Beta vs SPY", "type": "number", "chart_kind": None, "active": True}

r_m = data.assets.returns("SPY", timeframe="D1")   # aligned to the equity dates
both = returns.notna() & r_m.notna()
result(returns[both].cov(r_m[both]) / r_m[both].var(), "Beta vs SPY", fmt="ratio")
\end{codiceInterno}

Tre regole fanno il comportamento.

\begin{description}[leftmargin = 1.4em, style = nextline, itemsep = 4pt]
  \item[Il timeframe predefinito] È quello della strategia dietro la pagina,
    letto dal \cod{DESCRIPTOR} per albero sintattico, senza importare la
    strategia (\cod{strategy\_timeframes}). Per una strategia su più
    timeframe vale il più basso, cioè il più fine; per approcci e pagine
    senza strategia D1; in un ensemble il più basso fra le strategie dei suoi
    stream. Il codice ne sceglie un altro con \cod{timeframe="H1"}.
  \item[Il timeframe che manca] Se su disco non c'è, lo si costruisce da uno
    più fine che lo divide esattamente (H4 da H1 prima che da M5, D1 da M5),
    con le regole OHLCV del motore. Se non si può, \cod{DataError} elenca i
    timeframe che ci sono e propone \cod{timeframe=...}.
  \item[L'allineamento] \cod{series}, \cod{close} e \cod{returns} sono
    allineate alle date della metrica con la regola del motore: alla data
    $t$ l'ultimo valore noto a $t$, mai uno successivo. I dati intraday
    prima si chiudono al giorno, così la data $t$ riceve la chiusura di $t$.
\end{description}

\begin{nota}{Un dato che manca non è un errore del codice}
\cod{DataError} (sottoclasse di \cod{LookupError}) arriva alla pagina come
\enquote{Data not available}, non come errore. Studio salva comunque la
metrica, con un avviso: funzionerà quando il dato sarà caricato. La prova che
precede il salvataggio gira sullo stream con più dati, non su un'equity
finta: una metrica che legge \cod{data} si prova davvero.
\end{nota}

\subsection{Grafici a più serie}

Con \cod{chart\_kind: "figure"} un grafico ha più serie, punti con
etichetta, dispersioni, barre, aree e linee di riferimento, su un asse x di
date oppure di numeri --- quello che serve per la Security Market Line del
CAPM, dove sull'asse x c'è il beta. Il runtime produce una specifica JSON
(\cod{kind: "figure"}, \cod{traces}); \cod{plotCustomFigure} di
\file{static/js/charts.js} la disegna con Plotly.

\begin{codiceInterno}{python}{Un grafico con l'asse x numerico}{}
fig = chart("CAPM - Security Market Line", x="Beta", y="Annual return", y_format="pct")
fig.line([0.0, 2.0], [rf, rf + 2.0 * premium], name="SML", dash="dash")
fig.point(1.0, market, name="S&P 500")
fig.point(beta, holding, name="Holding", label="beta %.2f" % beta)
\end{codiceInterno}

\subsection{La pagina di uno stream}

Ogni metrica della pagina, predefinita o personale, numerica o grafica, si
nasconde e si rimette. Le predefinite sono gli elementi con \cod{data-sv} in
\file{templates/stream.html}; le personali le dispone
\file{static/js/stream\_metrics.js}: i numeri dopo quelli predefiniti, i
grafici in fondo alla pagina, dopo grafici e tabelle predefiniti, con
\voce{New chart metric} per ultima.

\begin{center}
\begin{tabular}{@{}L{44mm}L{94mm}@{}}
\toprule
\textbf{Pezzo} & \textbf{Che cosa fa} \\
\midrule
\rigaTenue
\cod{/api/stream-view/prefs} & la scelta, una per tutti gli stream:
  \cod{shown} (le chiavi accese) e \cod{known} (quelle esistenti quando la
  scelta è stata fatta). \cod{shown: null} = predefinite più personali
  attive \\
\cod{sv\_off(chiave)} & nel template: le predefinite nascoste partono
  nascoste, senza comparire per un istante \\
\rigaTenue
\cod{/api/metrics/compute} & con \cod{stream}, \cod{ids} e
  \cod{include\_inactive}: calcola solo le personali accese, sui risultati
  dello stream e con il suo benchmark \\
\cod{/studio?tab=metric} & con \cod{output=number} o \cod{output=chart}
  (\voce{New number metric}, \voce{New chart metric}): l'uscita scelta, un
  nome libero e, con \cod{from=}, il ritorno allo stream; con
  \cod{metric=<id>} riapre una metrica. \voce{Save} valida e salva in un
  colpo solo \\
\bottomrule
\end{tabular}
\end{center}

\begin{perche}{Perché \cod{known} oltre a \cod{shown}}
Una metrica creata dopo la scelta --- per esempio con \voce{New chart
metric} dalla pagina stessa --- non è stata scartata da nessuno. Senza
\cod{known} resterebbe nascosta finché qualcuno non va a cercarla
nell'elenco; con \cod{known} compare da sola, se è attiva.
\end{perche}

\subsection{La guardia: numpy e pandas filtrati}

Il codice di una metrica o di un indicatore può arrivare dalla Community.
numpy e pandas gli arrivavano così com'erano, e da un loro sottomodulo si
raggiungeva il sistema operativo senza scrivere un import
(\cod{pd.core.config\_init.os}); \cod{pd.read\_pickle} di un indirizzo
eseguiva codice arrivato dalla rete. \file{dashboard/sandbox\_guard.py} li
passa filtrati: funzioni e classi restano quelle vere, i sottomoduli solo
quelli di calcolo (\cod{np.linalg}, \cod{np.random}, \cod{pd.offsets},
\ldots); lettura e scrittura di file, \cod{eval}, \cod{query} e \cod{plot}
sono rifiutati già nel sorgente, anche chiamati per nome
(\cod{agg("to\_parquet")}). I dati del sistema si leggono con \cod{data}.

\begin{attenzione}{Uno strato, non un isolamento}
Come il validatore AST delle strategie (\file{adapters/ast\_validator.py}),
la guardia lavora nel processo dell'applicazione: chiude gli abusi diretti,
non un attaccante che costruisce i nomi a pezzi. La barriera vera resta
l'isolamento di processo. Il validatore delle strategie ha la stessa
apertura su numpy e pandas: una strategia può arrivare a \cod{os} per la
stessa strada.
\end{attenzione}
